\section{Background}
\label{sec:background}

\subsection{Prefix caching}
A decoder-only transformer computes, for each prompt token, key and value
vectors at every layer. A token's KV depends only on itself and the tokens
before it, so KV computed for one prompt is valid for any later prompt that
shares the same prefix~\cite{qin2025mooncake}. Engines store KV in fixed-size
blocks and key each block by a hash chained over its own tokens and its
prefix, $h_j = H(h_{j-1}, \mathrm{tokens}_j)$; a lookup returns the longest
matching chain from the first block~\cite{kwon2023vllm,addagada2026kvgov}. A
hit skips that part of prefill, cutting TTFT roughly in proportion to the
prefix length.

\subsection{\mc: a KVCache-centric architecture}
\mc~\cite{qin2025mooncake} separates prefill and decode onto different node
pools and turns the KVCache into a cluster-wide storage tier.

\paragraph{\mc Store.} KV blocks of 16--512 tokens (256 in the paper's
evaluation) live in a distributed pool built from the host DRAM and SSD of
every node. Each block's key covers the block and its prefix, so identical
content is deduplicated cluster-wide. Hot blocks are replicated across nodes to
avoid fetch congestion, and when the pool fills, blocks are evicted LRU unless in
use. A transfer engine moves blocks over multiple RDMA NICs per node; the paper
reports 87\,GB/s and 190\,GB/s on 4$\times$200 and 8$\times$400\,Gbps fabrics.
The released store adds per-object replica counts and soft and hard pins~\cite{mooncakerepo}.

\paragraph{Conductor.} A global scheduler places each request
(Algorithm~1 in~\cite{qin2025mooncake}). It hashes the prompt into block keys,
finds each prefill instance's prefix-match length, and estimates TTFT per
instance as transfer time for any remote blocks plus queueing plus prefill time,
the last from a regression fitted offline on request length and match length.
It picks the instance with the lowest estimate, and if the TTFT or TBT SLO
cannot be met, returns HTTP 429 instead of serving the request. Placement
decisions also trigger replication of hot blocks.

\paragraph{Why it matters here.} Three design choices that make \mc fast also
make it observable: (i) deduplication is keyed by content alone, with no
principal in the key; (ii) admission is a function of the \emph{global} match
length; and (iii) a block's physical location---local HBM, local DRAM, remote
DRAM, SSD---depends on other requests' access history. The paper evaluates a
single first-party service and does not discuss tenancy or
confidentiality~\cite{qin2025mooncake}.

\subsection{Multi-tenant serving}
LLMs are mostly sold as APIs to many mutually distrusting customers. Responses
expose TTFT and inter-token time (TBT) to the client, HTTP status codes, and
usage metadata that typically includes the number of cached prompt tokens,
which providers bill at a discount~\needcite{provider docs}. Engines offer
per-request \texttt{cache\_salt} to separate tenants' key spaces; NVIDIA Dynamo
uses the same salt to namespace KV-aware routing~\cite{vllmcachesalt,dynamonvext}.
Salting is only as strong as its plumbing: one 2026 vLLM bug dropped the salt on
tool continuations, restoring a cached-token oracle~\cite{cve2026105752}.
